% !TITLE 望岳
% !AUTHOR 杜甫
% !DYNASTY 唐
% !ORDER 18

\begin{poem}
岱宗夫如何？齐鲁青未了。
造化钟神秀，阴阳割昏晓。
荡胸生曾云，决眦入归鸟。
会当凌绝顶，一览众山小。
\end{poem}

\begin{appreciation}
这是杜甫现存最早的一首诗，写于他青年时代漫游齐鲁之时。那时的杜甫还没有经历安史之乱，还没有写出三吏三别，还是一个意气风发的年轻人。

岱宗夫如何？齐鲁青未了——岱宗就是泰山，五岳之首。夫如何是一个问句：泰山到底怎么样呢？答案是齐鲁青未了——在齐鲁大地上，泰山的青色延绵不绝，一眼望不到边。未了是未尽、无边的意思。这个问答式的开头，既写出了泰山的高大，也写出了诗人面对它时的惊叹。

造化钟神秀，阴阳割昏晓——大自然把神奇和秀美都集中在了泰山身上；山的南北两面，一面是阳光，一面是阴影，像被刀割开一样分明。钟是聚集的意思，割是分割的意思。这两个字都很有力量：造化偏爱泰山，所以把所有的美都给了它；泰山如此高大，所以能把阴阳两面割开。这是青年杜甫的笔触，充满了赞叹和敬畏。

荡胸生曾云，决眦入归鸟——山上的云气在胸中激荡，睁大眼睛看着归巢的飞鸟进入山林。决眦是睁裂眼角的意思，形容看得极认真、极用力。一个年轻人站在山下，仰望山顶，恨不得把眼睛睁到最大，把一切都看进去。这种对世界的贪婪的凝视，是青春独有的姿态。

会当凌绝顶，一览众山小——我一定要登上泰山的最高峰，到那时，所有的山都会显得渺小。这是全诗最著名的一句，也是青年杜甫精神的最好写照。他相信，只要努力，就能到达最高处；只要到达最高处，就能看到别人看不到的风景。这种自信，这种豪情，和后来那个沉郁顿挫的杜甫判若两人。但正是这个青年时代的杜甫，让我们看到了一个伟大诗人最初的模样：他对世界充满好奇，对未来充满信心，对自己充满期待。
\end{appreciation}
